\documentclass{article}
\usepackage{readmasters}
\begin{document}

\origpage{267}

\section*{Ueber die algebraischen Functionen einer und zweier Variabeln.}

\subsection*{Note 2${}^{1)}$. Von Max Noether in Heidelberg. (Mitgetheilt durch A.~Clebsch.)}

\section*{I.}

Man kann eine allgemeine Methode angeben, um eine gegebene Curve mit beliebigen singulären Punkten rational in eine andere zu transformiren, welche nur gewöhnliche vielfache Punkte enthält. Herr Cayley hat gezeigt${}^{2)}$, wie man die Puiseux'schen Reihenentwicklungen in einem singulären Punkte einer Curve benutzen kann, um die einem solchen Punkte in Bezug auf die

\textbf{1)} Vgl.\ Note 1 in diesen »Nachrichten«, vom 14.\ Juli 1869.

\textbf{2)} Quarterly Journal of Mathematics, 1866, vol.\ VII, p.\ 212.

\origpage{268}
Irrationalität der Curve äquivalente Anzahl von Doppel- und Rückkehrpunkten zu erhalten. Diese nämliche Bestimmung kann man mit Hülfe der erwähnten Transformation direct, ohne Voraussetzung der Pniseux'schen\ednote{Printed Pniseux'schen with an n instead of a u; an apparent misprint for Puiseux'schen.} Entwicklungen, erreichen.

Dazu genügt es, irgend eine rationale Transformation, bei welcher ein Fundamentalpunkt in den singulären Punkt $P$ der Curve $C$ gelegt wird, auf die Curve anzuwenden. Es wird dabei nur eine so allgemeine Lage der Transformationscurven gegen $C$ vorausgesetzt, dass die Jacobi'sche Curve der Transformation von den Tangenten von $C$ in $P$ nicht berührt wird.

Bei einer solchen Transformation löst sich die von dem vielfachen ($\nu$fachen) Punkte $P$ als solchem herrührende Singularität in der transformirten Curve $C'$ auf; und es bleiben demgemäss, dem Punkte $P$ entsprechend, auf $C'$ Punkte von niedrigerer Singularität, die zusammen, verbunden mit einem allgemeinen $\nu$fachen Punkte, äquivalent sind dem Punkte $P$ von $C$. Bei einer fortgesetzten Anwendung von Transformationen auf die singulären Punkte der so entstehenden Curven $C'$, $C'', \ldots$ erniedrigt sich somit die Singularität der $P$ entsprechenden Punkte immer mehr, bis endlich dem Punkte $P$ eine Reihe von einfachen Punkten der transformirten Curve entspricht.

Diese Auflösung der Singularität von $P$ in die von mehreren Punkten ist identisch mit der Trennung der Functionswerthe um den singulären Punkt in Klassen. Die Anzahl der getrennten einfachen Punkte, die zuletzt $P$ entsprechen, ist gleich der Anzahl der cyklischen

\origpage{269}
Systeme der Functionswerthe um den singulären Punkt${}^{1)}$.

Die Singularität von $C$ in $P$ zählt, in Bezug auf das Geschlecht von $C$, für $\frac{\nu(\nu-1)}{1.2}$ Doppelpunkte, plus der Anzahl der Doppelpunkte, welche in den dem Punkte $P$ entsprechenden singulären Punkten von $C'$ enthalten sind; und sie bestimmt sich somit durch eine Fortsetzung dieser Abzählung bei den successiven Transformationen. Auch die Anzahl der darunter enthaltenen Rückkehrpunkte (und zugleich die Anzahl der in einem der cyklischen Systeme enthaltenen Wurzeln) bestimmt sich aus der Art der Berührung von $C'$ mit der $P$ entsprechenden Fundamentalcurve der Transformation; eine Berührung $\mu^{\text{ter}}$ Ordnung eines Zweiges von $C'$ mit dieser Curve bedeutet ein cyklisches System von $\mu+1$ Wurzeln oder $\mu$ Rückkehrpunkte unter der Zahl der hiervon herrührenden Doppelpunkte.

Die einfachste und bequemste unter den hier anzuwendenden Transformationen ist die ebene quadratische, bei welcher den Geraden der Ebene Kegelschnitte durch 3 feste Punkte entsprechen.

So führt man die hyperelliptische Curve $2n^{\text{ter}}$ Ordnung mit singulärem $(2n-2)$fachem Punkte $P$ ($x_1 = x_2 = 0$)
\[ C = x_3^2 f_{n-1}(x_1, x_2) + f_{2n}(x_1, x_2) = 0, \]
wo, wie auch im Folgenden, der Index an $f$ die Ordnung der Functionen $f$ anzeigt, durch die quadratische Transformation

\textbf{1)} S.\ Puiseux, in Lionville's\ednote{Printed Lionville's with an n instead of a u; an apparent misprint for Liouville's.} Journal, t.\ 15 und 16.

\origpage{270}
\[ y_1 : y_2 : y_3 = x_2 x_3 : x_3 x_1 : x_1 x_2 \]
über in die Curve
\[ C' = y_1^2 y_2^2 f_{n-1}^2(y_2, y_1) + f_{2n}(y_2, y_1) \cdot y_3^2 = 0. \]\ednote{Printed with an exponent 2 on $f_{n-1}$; an apparent misprint for $f_{n-1}$.}

Dem singulären Punkte von $C$ entsprechen auf $C'$ die $n-1$ gewöhnlichen Doppelpunkte $y_3 = 0$, $f_{n-1}(y_2, y_1) = 0$, und $P$ ist
\[ \frac{(2n-2)(2n-3)}{1.2} + (n-1) \]
Doppelpunkten äquivalent, oder das Geschlecht von $C$ ist $n-1$. Es bilden sich hier zuerst $n-1$ Klassen, von denen jede sodann in zwei zerfällt; d.\ h.\ $C$ hat einen $(2n-2)$fachen Punkt mit $n-1$ getrennten Tangenten, in deren jeder sie einen Selbstberührungspunkt hat.

\section*{II.}

Die Methode dieser Transformation lässt sich auf algebraische Functionen zweier Variabeln ausdehnen und führt hier insbesondere zur Untersuchung singulärer Knotenpunkte einer Fläche und zur Bestimmung der Wirkung derselben auf das Flächengeschlecht, das bis jetzt erst bei vielfachen Curven und allgemeinen konischen Knotenpunkten der Fläche festgestellt worden ist${}^{1)}$.

Wenn man eine Raumtransformation${}^{2)}$ anwendet,

\textbf{1)} S.\ meine schon citirte Note vom 14.\ Juli 1869.

\textbf{2)} Ueber die hierzu geeignetsten Transformationen, die direct umkehrbaren, vgl.\ Cayley, Proc.\ of the London Math.\ Soc., vol.\ III, 1870.
Cremona, diese »Nachrichten« vom 4.\ Mai 1871, und Rendiconti del R.\ Istit.\ Lombardo, vom 5.\ Mai 1871, sowie meine gleichzeitig mit diesen beiden Noten erschienene Abhandlung in Math.\ Ann., Bd.\ 3, p.\ 547.

\origpage{271}
bei welcher ein Fundamentalpunkt in den singulären Punkt $P$ der Fläche $F$ gelegt wird, so entspricht dem Punkte $P$ auf der transformirten Fläche $F'$ eine Curve $C'$, durch welche $F'$ im Allgemeinen vielfach und singulär gehen wird, aber so, dass die Ordnung dieser Singularität niedriger ist, als die von $P$ auf $F$. Man kann nun entweder dieses Verfahren wiederholen, indem man $C'$ als Fundamentalcurve einer zweiten Transformation annimmt, wodurch man auf Curven von niedrigerer Singularität geführt wird; oder, was im Allgemeinen vortheilhafter ist, man bestimmt direct den Einfluss der vielfachen Curve von $F'$ auf das Flächengeschlecht dieser Fläche. Das Flächengeschlecht ist, wenn $F'$ von der Ordnung $m$ ist, gleich der Anzahl der Flächen $(m-4)^{\text{ter}}$ Ordnung, welche sich in den vielfachen Curven genau so zu verhalten haben, wie die Curven $(m-3)^{\text{ter}}$ Ordnung, welche das Geschlecht eines ebenen Querschnitts der Fläche bestimmen, in den Schnittpunkten dieses Querschnitts mit den vielfachen Curven. Man erschliesst daher aus den Reihenentwicklungen in einem ebenen Querschnitt in einem Punkte von $C'$, die denen in einem ebenen Querschnitt durch $P$ bei $F$ entsprechen und nur von niedrigerer Ordnung sind, das Verhalten der Flächen $(m-4)^{\text{ter}}$ Ordnung, und dadurch das Flächengeschlecht $p$ von $F$.

Diese Reihenentwicklungen werden, wenn man den ebenen Querschnitt unbestimmt lässt, in einzelnen Punkten der vielfachen Curven ungültig; aber für das Verhalten der Flächen $(m-4)^{\text{ter}}$ Ordnung ist die Untersuchung dieser Punkte,

\origpage{272}
die übrigens selbst nach der angegebenen Methode geschehen könnte, überflüssig.

Am bequemsten zur Transformation ist die quadratische:
\[
\begin{aligned}
x_1 : x_2 : x_3 : x_4 &= y_1 y_4 : y_2 y_4 : y_3 y_4 : \varphi_2(y_1, y_2, y_3) \\
y_1 : y_2 : y_3 : y_4 &= x_1 x_4 : x_2 x_4 : x_3 x_4 : \varphi_2(x_1, x_2, x_3)
\end{aligned}
\tag{1}
\]
bei welcher dem Puncte\ednote{Printed Puncte with a c, alongside Punkte with a k elsewhere in the text.} $P$ ($x_1 = x_2 = x_3 = 0$) die Ebene $E'$ ($y_4 = 0$), welche noch den Fundamentalkegelschnitt $K'$ ($y_4 = 0$, $\varphi_2(y) = 0$) enthält, entspricht, und einer Berührungsebene $A(x_1, x_2, x_3) = 0$ von $F$ im Punkte $P$ eine Gerade $A(y_1, y_2, y_3) = 0$, $y_4 = 0$ von $F'$.

Ich erlaube mir, durch einige Beispiele die Anwendbarkeit der Methode zu zeigen.

a) Wenn die Singularität eines $\mu$fachen Punktes $P$ von $F$ darin besteht, dass der Tangentenkegel $\mu^{\text{ter}}$ Ordnung in $P$ eine $\nu$fache Kante besitzt, so erhält die nach 1) entsprechende Fläche $F'$ die Singularität, dass sie von der Ebene $E'$ in einem Punkte in der $(\nu-1)^{\text{ten}}$ Ordnung berührt wird, was auf das Geschlecht von $F'$ ohne Einfluss ist. Daher reducirt auch der specielle $\mu$fache Punkt $P$ das Geschlecht $p$ von $F$ nur um $\frac{1}{6}\mu(\mu-1)(\mu-2)$.

b) Sei $F$ wieder eine Fläche mit $\mu$fachen singulären Knotenpunkte $P$:
\[
\begin{aligned}
F = x_4^{n-\mu} f_\mu(x_1, x_2, x_3) &+ x_4^{n-\mu-1} f_{\mu+1}(x) \\
&+ \ldots + f_n(x) = 0,
\end{aligned}
\tag{2}
\]
und sei insbesondere $x_1 = x_2 = 0$ eine $(\nu-1)$fache

\origpage{273}
Kante für $f_{\mu+i}(x_1, x_2, x_3) = 0$, von $i = 0$ bis $i = \nu - 1$, für $\mu + \nu - 1 \leqq n$. Dann wird $y_1 = y_2 = y_4 = 0$ ein $\nu$facher Punkt von $F'$, und das Geschlecht $p$ von $F$ erniedrigt sich durch $P$ um
\[ \frac{1}{6}\mu(\mu-1)(\mu-2) + \frac{1}{6}\nu(\nu-1)(\nu-2). \]

So wird, für $n = 5$, $\mu = 3$, $\nu = 3$, die Fläche $F_5$, wenn sie noch eine Doppelgerade besitzt, auf der Ebene abbildbar.

c) Sei bei der Fläche 2) die Gerade $x_1 = x_2 = 0$ eine $\nu$fache Kante von $f_\mu(x) = 0$, $f_{\mu+1}(x) = 0, \ldots$ und $f_n(x) = 0$, für $\nu \leqq \mu$. Bei der speciellen Transformation 1) erhält dann $F'$ ein ähnliches System mit $n$fachem Punkte $P'$ ($y_1 = y_2 = y_3 = 0$) und $\nu$facher Geraden $y_1 = y_2 = 0$. Und man folgert aus der Gleichsetzung der Reduction auf $p$ für beide vielfache Gerade den Satz, dass, wenn eine Fläche einen $\mu$fachen Punkt enthält, der für sich allein die Reduction $M = \frac{1}{6}\mu(\mu-1)(\mu-2)$, und eine von $P$ ausgehende $\nu$fache Curve, welche für sich die Reduction $N$ auf das Geschlecht der Fläche hervorbringen würde, die gesammte Reduction durch Punkt und Curve
\[ M + N - \frac{\nu(\nu-1)(3\mu-2\nu-2)}{6}, \quad \text{für } \nu \leqq \mu, \]
beträgt.

d) $F$ besitze einen uniplanaren Knotenpunkt $\mu^{\text{ter}}$ Ordnung, aber der Art singulär, dass $F$ die Gleichung hat:

\origpage{274}
\[
\begin{aligned}
F = x_4^{n-\mu} x_1^\mu &+ x_4^{n-\mu-1} x_1^{\mu-1} f_2(x_1, x_2, x_3) \ldots \\
&+ x_4^{n-2\mu+1} x_1 f_{2\mu-2}(x) + x_4^{n-2\mu} f_{2\mu}(x) \ldots \\
&+ f_n(x) = 0, \qquad\qquad \text{für } 2\mu-1 > n.
\end{aligned}
\]

Dann erhält $F'$ nach 1) als Singularität noch eine $\mu$fache Gerade $y_1 = y_4 = 0$, durch welche sich das Geschlecht bei Berücksichtigung${}^{1)}$ der beiden Schnittpunkte dieser Geraden mit $K'$, um $\frac{1}{2}\mu(\mu-1)^2$ erniedrigt. Der singuläre uniplanare Knotenpunkt von $F$ reducirt daher das Geschlecht $p$ um
\[ \frac{1}{6}\mu(\mu-1)(\mu-2) + \frac{1}{2}\mu(\mu-1)^2 = \frac{1}{6}\mu(\mu-1)(4\mu-5). \]

Für $\mu = 2$ ergiebt sich speciell, dass ein uniplanarer Doppelpunkt, von der Art, dass jede durch denselben gehende Ebene die Fläche in einer Curve mit zwei zusammenfallenden Doppelpunkten schneidet, also ein Selbstberührungspunkte\ednote{Printed Selbstberührungspunkte with a trailing e; an apparent misprint for the singular Selbstberührungspunkt.} der Fläche, das Geschlecht derselben um 1 erniedrigt, und eine Fläche $4^{\text{ter}}$ Ordnung
\[ x_4^2 x_1^2 + 2 x_4 x_1 f_2(x_1, x_2, x_3) + f_4(x_1, x_2, x_3) = 0 \]
muss sich auf der Ebene abbilden lassen. In der That führt diese Fläche auch auf eine Doppelebene mit Uebergangscurve $4^{\text{ter}}$ Ordnung${}^{2)}$.

\textbf{1)} S.\ die »postulation«-Formel in dem schon citirten Cayley'schen Memoire über die Raumtransformationen, pag.\ 180.

\textbf{2)} S.\ Clebsch, Math.\ Ann., Bd.\ 3, pag.\ 51.

\origpage{275}
Die ebenen Schnitte der Fläche bilden sich durch Curven $6^{\text{ter}}$ Ordnung ($1^2$, $2^2, \ldots 7^2$, $8$, $9$, $10$, $11$) ab, wobei die Punkte $8$, $9$, $10$, $11$ von einer Curve $3^{\text{ter}}$ Ordnung ($1$, $2, \ldots 7$) aus einer der Curven $6^{\text{ter}}$ Ordnung ausgeschnitten werden. Das System der 28 durch den uniplanaren Punkt gehenden noch doppelt berührenden Ebene\ednote{Printed Ebene in the singular; an apparent misprint for the plural Ebenen.} ist identisch mit dem der Doppeltangenten einer Curve $4^{\text{ter}}$ Ordnung${}^{1)}$.

\section*{III.}

Ich mache noch eine Anwendung dieser Theorie auf die Untersuchung der Bedingung, unter welcher eine geometrische Doppelebene auf einer Kegelfläche abbildbar wird. Von Abbildungen auf der einfachen Ebene hat Herr Clebsch in seiner schon citirten Abhandlung, Math.\ Ann.\ Bd.\ 3, Beispiele gegeben.

Sei die »Uebergangscurve« der Doppelebene eine Curve der Ordnung $2m$, $\varOmega_{2m}(x_1, x_2, x_3) = 0$. Ich betrachte die Fläche
\[ F = x_4^2 f_{m-1}^2(x_1, x_2, x_3) - \varOmega_{2m}(x_1, x_2, x_3) = 0, \]
die mit der Doppelebene zugleich auf einer Kegelfläche abbildbar wird. Ihre Singularität liegt in dem $(2m-2)$fachen Punkte $P$ ($y_1 = y_2 = y_3 = 0$). Die Transformation von 1) II. liefert
\[ F' = \varphi_2^2 f_{m-1}^2(y_1, y_2, y_3) - y_4^2 \varOmega_{2m}(y_1, y_2, y_3) = 0, \]

\textbf{1)} Ein specieller Fall dieser Fläche findet sich bei Herrn Cremona, diese Nachrichten vom 3.\ Mai 1871, p.\ 714, erwähnt.

\origpage{276}
eine Fläche der Ordnung $2m+2$, mit allgemeinem $2m$fachen Punkte $P'$ ($y_1 = y_2 = y_3 = 0$), mit dem Doppelkegelschnitt $K'$ und mit einer in dessen Ebene liegenden allgemeinen Doppelcurve $(m-1)^{\text{ter}}$ Ordnung ($y_4 = 0$, $f_{m-1}(y) = 0$). Das Geschlecht dieser Fläche ist $\frac{1}{2}(m-1)(m-2)$.

»Das Flächengeschlecht der Flächen, welche auf eine Doppelebene mit allgemeiner Uebergangscurve führen, ist $\frac{1}{2}(m-1)(m-2)$«.

Wenn $\varOmega_{2m}(x) = 0$ einen $\nu$fachen Punkt in $x_2 = x_3 = 0$ besitzt, so erhält $F$ einen singulären uniplanaren Doppelpunkt im Punkte $\varPi$ ($x_2 = x_3 = x_4 = 0$). Man macht daher eine Transformation
\[ x_2 : x_3 : x_4 : x_1 = \xi_2 \xi_1 : \xi_3 \xi_1 : \xi_4 \xi_1 : \psi_2(\xi_2, \xi_3, \xi_4) \]
und erhält eine Fläche $\varPhi$ mit singulärer Geraden $G$ ($\xi_1 = \xi_4 = 0$). Für ein gerades $\nu$, $\nu = 2\varrho$, ergeben sich sodann in einem Punkte von $G$ zwei getrennte Reihenentwicklungen:
\[
\begin{gathered}
\xi_4 = \varkappa \xi_1^{\varrho-1} + \ldots \\
\xi_4 = -\varkappa \xi_1^{\varrho-1} + \ldots,
\end{gathered}
\]
und die Fläche $\varPhi$ hat zwei Schalen, welche sich $(\varrho-1)$punktig in jedem Punkte der Geraden $G$ treffen, in der singulären Berührungsebene $\xi_4 = 0$. Man leitet hieraus ab, dass die Fläche $(2m-4)^{\text{ter}}$ Ordnung, deren Anzahl das Geschlecht $p$ der Fläche $F$ bestimmt, durch den Punkt $\varPi$ derart hindurchgehen müssen, dass sie daselbst die Ebene

\origpage{277}
$x_4 = 0$ in der $(\varrho-2)^{\text{ten}}$ Ordnung berühren, was $\frac{\varrho(\varrho-1)}{1.2}$ Bedingungen darstellt.

Wenn $\nu$ ungerade${}^{1)}$, $\nu = 2\varrho+1$, so erhält man in einem Punkte von $G$ ein cyklisches System von 2 Wurzeln und die Reihenentwicklung für diese beiden Wurzeln:
\[ \xi_4 = \pm \varkappa \xi_1^{\frac{2\varrho-1}{2}}. \]
Die Fläche $\varPhi$ hat dann in $\xi_1 = \xi_4 = 0$ zwei Schalen, die sich in $\varrho-1$ zusammenfallenden Geraden schneiden, und die singuläre Doppellinie ist äquivalent mit $(\varrho-2)$ Doppelgeraden und einer Rückkehrgeraden. Auch hier folgt, dass jene Flächen $(2m-4)^{\text{ten}}$ Ordnung die Ebene $x_4 = 0$ in $P$\ednote{Printed $P$; in the preceding text on p.~276 this point was denoted $\varPi$.} in der $(\varrho-2)^{\text{ter}}$ Ordnung berühren müssen, was wieder $\frac{\varrho(\varrho-1)}{1.2}$ Bedingungen giebt.

»Das Geschlecht einer Fläche, welche auf eine Doppelebene mit Uebergangscurve der Ordnung $2m$ führt, die $\alpha_i$ $2i$fache oder $(2i+1)$fache Punkte besitzt, ist $=\frac{1}{2}(m-1)(m-2) - \Sigma_i \frac{1}{2} i(i-1)\alpha_i$.

Dass dieser Ausdruck verschwinde, ist die Bedingung für die Abbildbarkeit der Doppelebene auf der einfachen Ebene. Ferner muss dieser Ausdruck $=-q$ werden, wenn die Doppelebene auf einem Kegel abbildbar sein soll, dessen ebene Querschnitte das Geschlecht $q$ haben.«

\textbf{1)} Hiernach ist eine Bemerkung von Zeuthen, Math.\ Ann.\ Bd.\ 3, pag.\ 324 zu corrigiren.

\origpage{278}
Diese letztere Bemerkung folgt aus einem Satze des Herrn Cayley${}^{1)}$, dass das Flächengeschlecht der Flächen, für welche die numerische Formel dasselbe $\leqq 0$ giebt, gleich wird dem »Curvengeschlecht«${}^{2)}$, negativ genommen.

Daher muss sich z.\ B.\ eine Doppelebene mit Uebergangscurve $2m^{\text{ter}}$ Ordnung, welche einen $(2m-2)$fachen Punkt besitzt, auf der Ebene abbilden lassen. In der That führen auch die Flächen $(m+1)^{\text{ter}}$ Ordnung mit $(m-1)$facher Geraden, auf eine solche Doppelebene. Die Abbildung derselben geschieht genau nach dem von Herrn Clebsch für $m = 3$ angegebenen Verfahren${}^{3)}$. Den Geraden der Doppelebene kann man auf der einfachen Ebene Curven der Ordnung $n+1$, mit $(n-1)$fachem, und $4n-2$ einfachen festen Punkten, entsprechen lassen, und die Zahl solcher Abbildungen beträgt $2^{4n-4}$.

Eine Fläche $4^{\text{ter}}$ Ordnung mit 2 uniplanaren Doppelpunkten, die Selbstberührungspunkte sind, führt auf eine Uebergangscurve $4^{\text{ter}}$ Ordnung mit $4$fachem Punkte. Daher lässt sich diese Fläche auf einem Kegel dritter Ordnung abbilden${}^{4)}$.

\textbf{1)} Nach einer Privatmittheilung von Herrn Cayley.

\textbf{2)} S.\ die erste Note vom 14.\ Juli 1869.

\textbf{3)} Math.\ Ann., Bd.\ 3, pag.\ 64.

\textbf{4)} Ebd.\ pag.\ 558.

\emph{Heidelberg}, den 26.\ Mai 1871.

\end{document}
